\chapter{PHÁT TRIỂN HỆ THỐNG TÌM KIẾM}\label{chap:timkiem}

Tìm kiếm kết hợp hai tín hiệu bổ sung cho nhau. BM25 bắt đúng thuật ngữ hiếm (tên mô hình, tên bộ dữ liệu), còn
embedding đa ngôn ngữ bắt được cách diễn đạt khác và câu hỏi tiếng Việt. Hai danh sách kết quả được gộp bằng
Reciprocal Rank Fusion (RRF)~\cite{rrf}, rồi quy về cấp bài báo.

\section{Chia đoạn}

\begin{itemize}
\item \textbf{Đoạn tóm tắt} (chỉ số 0): tiêu đề + tóm tắt. Bài chỉ có metadata (không có PDF) vẫn tìm được nhờ
  đoạn này.
\item \textbf{Đoạn nội dung:} cửa sổ trượt 220 từ, chồng lấn 40 từ, không vượt qua ranh giới mục; mỗi đoạn mang
  số trang và tên mục để trích dẫn chính xác.
\item Phần References bị cắt khỏi nội dung cần chia đoạn, nhưng các mã arXiv trong đó được giữ lại làm cạnh CITES.
\end{itemize}

\section{Embedding}

\begin{table}[H]
\centering\small\setstretch{1.15}
\begin{tabularx}{\linewidth}{@{}P{0.22\linewidth}L@{}}
\toprule
\textbf{Thuộc tính} & \textbf{Giá trị}\\
\midrule
Model mặc định & \code{intfloat/multilingual-e5-base}~\cite{me5}: khoảng 278 triệu tham số, vector 768 chiều, hỗ trợ tiếng Việt\\
Tiền tố & \code{query:} cho câu hỏi, \code{passage:} cho văn bản; họ BGE dùng câu chỉ dẫn riêng, tự nhận theo tên model\\
Chuẩn hóa & Chuẩn L2, nên tích vô hướng giữa hai vector chính là cosine\\
Thiết bị & Tự chọn CUDA $\to$ MPS $\to$ CPU; trên GPU dùng số thực 16 bit và lô lớn gấp 4 lần\\
Chạy độc lập & \code{python -m scholarlens.embedding --device cuda}, chỉ xử lý các đoạn có \code{embedded = 0}\\
\bottomrule
\end{tabularx}
\caption{Cấu hình embedding}\label{tab:embedding}
\end{table}

\section{Tìm kiếm lai và gộp hạng}

\paragraph{BM25.} Mỗi đoạn văn $d$ được chấm điểm theo các từ $t$ của câu hỏi $q$:
\begin{equation}
\BM(q,d) = \sum_{t \in q} \mathrm{IDF}(t)\,\frac{f(t,d)\,(k_1+1)}{f(t,d) + k_1\bigl(1 - b + b\,\frac{|d|}{\mathrm{avgdl}}\bigr)},
\label{eq:bm25}
\end{equation}
trong đó $f(t,d)$ là số lần từ $t$ xuất hiện trong đoạn (điểm tăng chậm dần khi từ lặp lại nhiều),
$\mathrm{IDF}(t)$ cho từ hiếm điểm cao hơn từ phổ biến, $|d|/\mathrm{avgdl}$ phạt nhẹ đoạn dài hơn trung bình;
$k_1 = 1{,}5$, $b = 0{,}75$ (mặc định của rank-bm25). Văn bản được tách từ đơn giản: chữ thường, bỏ từ dừng.

\paragraph{Tìm kiếm dense.} Câu hỏi được embedding (tiền tố \code{query:}), ChromaDB trả về 50 đoạn có cosine cao
nhất, có thể kèm điều kiện lọc theo bài hoặc theo năm.

\paragraph{Gộp bằng RRF.} Hai danh sách top-50 được gộp chỉ dựa trên thứ hạng:
\begin{equation}
\RRF(d) = \sum_{r \in \{\mathrm{dense},\,\BM\}} \frac{1}{k + \rank_r(d)}, \qquad k = 60.
\label{eq:rrf}
\end{equation}
Đoạn đứng cao ở cả hai danh sách được cộng điểm từ cả hai. RRF chỉ dùng thứ hạng nên không cần chuẩn hóa thang
điểm giữa BM25 (không có giới hạn trên) và cosine (từ $-1$ đến $1$).

\paragraph{Quy về cấp bài báo.} Các đoạn được nhóm theo bài; điểm của bài bằng điểm đoạn cao nhất cộng
$0{,}3 \times$ tổng điểm các đoạn còn lại của bài đó. Bài có nhiều đoạn khớp được thưởng nhẹ, nhưng một đoạn rất
khớp vẫn quyết định thứ hạng.

\section{Truy cập tìm kiếm}

\begin{table}[H]
\centering\small\setstretch{1.15}
\begin{tabularx}{\linewidth}{@{}P{0.38\linewidth}L@{}}
\toprule
\textbf{Điểm truy cập} & \textbf{Mô tả}\\
\midrule
Khung chat (chế độ Research) & Mọi câu hỏi chuyên môn đều chạy tìm kiếm lai; các bài tìm được hiện thành danh sách nguồn dưới câu trả lời\\
\code{GET /api/search?q=...\&mode=hybrid} & API trả về danh sách bài, điểm, đoạn trích và thời gian xử lý; \code{mode} nhận \code{hybrid}, \code{dense} hoặc \code{bm25} để so sánh ba phương pháp\\
\code{GET /api/papers?q=...} & Lọc nhanh kho theo tiêu đề, tác giả hoặc mã bài\\
\bottomrule
\end{tabularx}
\caption{Các điểm truy cập tìm kiếm}\label{tab:truycap}
\end{table}
